\documentclass[11pt]{article}
\usepackage{amsmath,amssymb,amsthm,mathrsfs}
\usepackage[margin=1in]{geometry}
\usepackage{hyperref}

\newtheorem{theorem}{Theorem}
\newtheorem{proposition}{Proposition}
\newtheorem{assumption}{Assumption}

\title{Finite-Mode Bohmian--Hartree Dynamics with Oscillator Memory}
\author{}
\date{}

\begin{document}
\maketitle

\begin{abstract}
We study a coupled Schr\"odinger--ODE system in which a nonrelativistic wavefunction
interacts through a finite family of Hartree-type nonlocal terms whose strengths are
modulated by classical harmonic modes. The model admits an action principle and a conserved
total energy. Bohmian trajectories are defined by the usual guidance law, and $|\Psi|^2$-equivariance
holds wherever the associated velocity field generates a flow avoiding nodes. The framework is
best interpreted as an effective quantum--classical closure for a system interacting with a small
set of environmental or feedback degrees of freedom (finite memory), not as a fundamental modification
of quantum mechanics.
\end{abstract}

\section{State space and notation}
Let $d\in\{1,2,3\}$ and $\Psi:\mathbb{R}^d\times[0,\infty)\to\mathbb{C}$ with $\|\Psi(t)\|_{L^2}=1$.
Set $\rho(x,t)=|\Psi(x,t)|^2$. Let $\mathcal{I}=\{1,\dots,M\}$ be a finite index set (optionally
$\mathcal{I}=\{p_1,\dots,p_M\}$ the first $M$ primes, used only as labels).

For each $k\in\mathcal{I}$ introduce a classical oscillator coordinate $q_k(t)\in\mathbb{R}$ and
momentum $\pi_k(t)=\dot q_k(t)$. Fix $\omega_k>0$, couplings $g_k\in\mathbb{R}$, and kernels
$K_k\in L^1(\mathbb{R}^d)$ that are even: $K_k(x)=K_k(-x)$.

Define the Hartree functionals
\[
H_k[\rho]\;:=\;\frac12\int_{\mathbb{R}^d}\int_{\mathbb{R}^d}\rho(x)\,K_k(x-y)\,\rho(y)\,dx\,dy
\qquad (k\in\mathcal{I}).
\]

\section{Action principle and equations of motion}
Let $\hat H_0=-\frac{\hbar^2}{2m}\Delta + V_{\mathrm{ext}}(x)$ with real $V_{\mathrm{ext}}$.
Consider the action
\begin{align*}
\mathcal{S}[\Psi,\overline{\Psi},q]
=\int dt\Big(
&\ \mathrm{Re}\,\langle \Psi,\, i\hbar\,\partial_t\Psi\rangle
-\langle \Psi,\hat H_0\Psi\rangle
-\sum_{k\in\mathcal{I}} g_k\,q_k(t)\,H_k[\rho(\cdot,t)]
\\
&\ +\sum_{k\in\mathcal{I}}\Big(\frac12\dot q_k(t)^2-\frac12\omega_k^2 q_k(t)^2\Big)
\Big),
\end{align*}
with $\langle f,g\rangle=\int_{\mathbb{R}^d}\overline{f}g\,dx$. Stationarity yields:

\paragraph{Schr\"odinger equation (mode-modulated Hartree interaction).}
\begin{equation}
i\hbar\,\partial_t\Psi
=\hat H_0\Psi
+\Big(\sum_{k\in\mathcal{I}} g_k\,q_k(t)\,(K_k*\rho)(x,t)\Big)\Psi,
\qquad \rho=|\Psi|^2.
\tag{S}
\end{equation}

\paragraph{Oscillator equations (finite memory).}
\begin{equation}
\dot q_k=\pi_k,\qquad
\dot \pi_k=-\omega_k^2 q_k - g_k\,H_k[\rho(\cdot,t)].
\tag{O}
\end{equation}

\section{Conserved total energy}
Define
\begin{align*}
\mathcal{E}(t)=
\int_{\mathbb{R}^d}\Big(\frac{\hbar^2}{2m}|\nabla\Psi|^2 + V_{\mathrm{ext}}|\Psi|^2\Big)\,dx
+\sum_{k\in\mathcal{I}}\Big(\frac12\pi_k^2+\frac12\omega_k^2 q_k^2\Big)
+\sum_{k\in\mathcal{I}} g_k\,q_k\,H_k[\rho].
\end{align*}

\begin{proposition}[Energy conservation]
For sufficiently regular solutions of \textup{(S)}--\textup{(O)}, $\frac{d}{dt}\mathcal{E}(t)=0$.
\end{proposition}

\section{Bohmian trajectories and equivariance}
Define current $j=\frac{\hbar}{m}\mathrm{Im}(\overline{\Psi}\nabla\Psi)$ and, where $\Psi\neq0$,
\[
v^\Psi(x,t)=\frac{j(x,t)}{\rho(x,t)}
=\frac{\hbar}{m}\,\mathrm{Im}\Big(\frac{\nabla\Psi}{\Psi}\Big)(x,t).
\]
Trajectories solve
\begin{equation}
\dot X_t=v^\Psi(X_t,t).
\tag{B}
\end{equation}

\begin{proposition}[Continuity equation]
If $\Psi$ is a sufficiently regular solution of \textup{(S)}, then
$\partial_t\rho+\nabla\cdot(\rho v^\Psi)=0$.
\end{proposition}

\begin{theorem}[$|\Psi|^2$-equivariance (conditional on flow)]
Assume that for $t\in[0,T]$ the velocity field $v^\Psi(\cdot,t)$ generates a measurable flow
on $\mathbb{R}^d\setminus Z_t$ with $Z_t=\{x:\Psi(x,t)=0\}$, and that trajectories avoid $Z_t$.
If $X_0$ has density $\rho(\cdot,0)$, then $X_t$ has density $\rho(\cdot,t)$ for all $t\in[0,T]$.
\end{theorem}

\noindent
In the linear case, global existence of Bohmian trajectories for typical initial configurations
despite nodal singularities is established in rigorous work (e.g.\ Berndl et al.).%
\footnote{See K.\ Berndl et al., ``Global existence and uniqueness of Bohmian trajectories'', arXiv:quant-ph/9509009.}

\section{Well-posedness: minimal, checkable hypotheses}
The PDE \textup{(S)} is a semilinear/nonlinear Schr\"odinger equation with a Hartree-type nonlinearity
and time-dependent coefficients $q_k(t)$. The ODE \textup{(O)} is smooth and finite-dimensional.

\begin{assumption}[Baseline regularity]
$V_{\mathrm{ext}}\in W^{2,\infty}(\mathbb{R}^d)$, $K_k\in L^1(\mathbb{R}^d)$ even, and
$\Psi_0\in H^2(\mathbb{R}^d)$ with $\|\Psi_0\|_{L^2}=1$, and $q_k(0),\pi_k(0)\in\mathbb{R}$.
\end{assumption}

\begin{theorem}[Local well-posedness (standard)]
Under the baseline regularity, there exists $T>0$ and a unique solution with
$\Psi\in C([0,T];H^2)$ and $(q,\pi)\in C^1([0,T];\mathbb{R}^{2M})$.
Moreover $\|\Psi(t)\|_{L^2}=1$ for $t\in[0,T]$.
\end{theorem}

\noindent
Sharper global results depend on dimension and kernel class (e.g.\ additional decay/smoothness or
energy-subcritical structure); the Hartree literature provides many such theorems.%
\footnote{For examples of global well-posedness results in Hartree-type equations in various function spaces,
see e.g.\ Y.\ Cho \& K.\ Nakanishi (survey notes) and recent work such as R.\ Manna (2022) for Hartree-type global results.}

\section{Topology on multiply connected configuration spaces}
Let the configuration space be a manifold $M$ (e.g.\ $\mathbb{R}^d\setminus\Omega$ with obstacles).
On a loop $\gamma$ avoiding nodes, write locally $\Psi=|\Psi|e^{iS/\hbar}$ and $v=\nabla S/m$.
Then the circulation
\[
\Gamma(t)=\oint_\gamma v(\cdot,t)\cdot d\ell
\]
is quantized in units $2\pi\hbar/m$ and invariant as long as nodes do not cross $\gamma$.

\section{Interpretation and relation to existing models}
The finite oscillator bank is an effective ``environment'' or feedback layer.
In the limit of many modes with suitable spectral density one recovers familiar oscillator-bath
phenomenology (Caldeira--Leggett / quantum Brownian motion).%
\footnote{See the overview ``Caldeira--Leggett model'', Scholarpedia.}
The structure is also analogous to mean-field and self-consistent Hartree models; the
Schr\"odinger--Newton equation is a well-known special case where $K(x)= -Gm^2/|x|$ in $d=3$.%
\footnote{For discussion and references, see e.g.\ the Schr\"odinger--Newton literature surveyed in arXiv:0803.4488.}
It is conceptually close to electron--phonon coupling models (polaron physics) where a small set of
collective modes modulates effective interactions.%
\footnote{For modern context and historical pointers to Fr\"ohlich-type coupling, see F.\ Giustino, Rev.\ Mod.\ Phys.\ 89, 015003 (2017).}

\section{Scope and empirical constraints}
As a universal fundamental replacement for linear quantum mechanics, nonlinear Schr\"odinger
dynamics is strongly constrained both conceptually and experimentally.
Weinberg's nonlinear framework motivated precision tests with null results and tight bounds.%
\footnote{S.\ Weinberg, Ann.\ Phys.\ 194 (1989); S.\ Weinberg, Phys.\ Rev.\ Lett.\ 62, 485 (1989);
J.\ J.\ Bollinger et al., Phys.\ Rev.\ Lett.\ 63, 1031 (1989).}
In addition, generic nonlinear modifications can enable superluminal signaling when combined with
entanglement (Polchinski; and re-derivations of Gisin-type arguments).%
\footnote{J.\ Polchinski, Phys.\ Rev.\ Lett.\ 66, 397 (1991); A.\ Bassi \& K.\ Hejazi, arXiv:1411.1768.}
Accordingly, the intended use here is as an \emph{effective} model for restricted regimes
(semiclassical feedback, coarse-grained environment, finite-memory closure), not a universal law.

\section{A minimal toy instance (computable baseline)}
In $d=1$ with a single mode and kernel $K(x)=e^{-|x|/\ell}$, $\ell>0$, the model reduces to
\[
i\hbar\Psi_t=-\frac{\hbar^2}{2m}\Psi_{xx}+V_{\mathrm{ext}}\Psi+g\,q(t)\,(K*|\Psi|^2)\Psi,\qquad
\ddot q+\omega^2 q=-g\,H[\rho],
\]
with Bohmian trajectories $\dot X_t = \frac{\hbar}{m}\mathrm{Im}(\Psi_x/\Psi)(X_t,t)$.
This serves as a direct numerical testbed (split-step FFT for $\Psi$ plus a standard ODE integrator for $q$).

\section{Optional indexing choice}
If one uses prime labels, it is purely a dictionary choice for incommensurate mode parameters
(e.g.\ $\omega_{p}\propto\log p$) and carries no number-theoretic physical claim.

\bibliographystyle{plain}
\begin{thebibliography}{9}
\bibitem{WeinbergAnnPhys1989} S.~Weinberg, \emph{Testing quantum mechanics}, Ann.\ Phys.\ \textbf{194} (1989).
\bibitem{WeinbergPRL1989} S.~Weinberg, \emph{Precision tests of quantum mechanics}, Phys.\ Rev.\ Lett.\ \textbf{62}, 485 (1989).
\bibitem{BollingerPRL1989} J.~J.~Bollinger et al., \emph{Test of the linearity of quantum mechanics by rf spectroscopy of trapped ions}, Phys.\ Rev.\ Lett.\ \textbf{63}, 1031 (1989).
\bibitem{Polchinski1991} J.~Polchinski, \emph{Weinberg's nonlinear quantum mechanics and the Einstein--Podolsky--Rosen paradox}, Phys.\ Rev.\ Lett.\ \textbf{66}, 397 (1991).
\bibitem{BassiHejazi2014} A.~Bassi and K.~Hejazi, \emph{No-faster-than-light-signaling implies linear evolutions}, arXiv:1411.1768.
\bibitem{Berndl1995} K.~Berndl et al., \emph{Global existence and uniqueness of Bohmian trajectories}, arXiv:quant-ph/9509009.
\bibitem{CLScholarpedia} \emph{Caldeira--Leggett model}, Scholarpedia.
\bibitem{GiustinoRMP2017} F.~Giustino, \emph{Electron-phonon interactions from first principles}, Rev.\ Mod.\ Phys.\ \textbf{89}, 015003 (2017).
\bibitem{SNReview2008} See e.g.\ discussion in \emph{Schr\"odinger--Newton equation as a possible generator of quantum state reduction}, arXiv:0803.4488.
\end{thebibliography}

\end{document}
